% Generated draft: requires mathematical and source review.
\documentclass[11pt]{article}
\usepackage[margin=25mm]{geometry}
\usepackage{fontspec}
\setmainfont{Latin Modern Roman}
\usepackage{amsmath,amssymb,mathtools,unicode-math}
\setmathfont{Latin Modern Math}
\usepackage{enumitem,needspace,xurl,hyperref}
\setlength{\parindent}{0pt}
\setlength{\parskip}{5pt}
\setlength{\emergencystretch}{3em}
\newcommand{\N}{\mathbb{N}}
\newcommand{\Z}{\mathbb{Z}}
\newcommand{\Q}{\mathbb{Q}}
\newcommand{\R}{\mathbb{R}}
\newcommand{\C}{\mathbb{C}}
\newcommand{\F}{\mathbb{F}}
\newcommand{\A}{\mathbb{A}}
\newcommand{\PP}{\mathbb{P}}
\newcommand{\E}{\mathbb{E}}
\newcommand{\eps}{\varepsilon}
\newcommand{\dd}{\,\mathrm{d}}
\begin{document}
\section*{UnsolvedMath ID: 30000708}
\textbf{Source problem code:} catalogue entry
\subsection{Definitions and mathematical statement}
The source asks for a Cram\'er--Wold-type uniqueness theorem beyond finite positive measures, allowing signed measures with singular or infinite mass near the origin. For a finite measure, halfspace data may be written
\[
H(u,t)=\{x\in\R^d:\langle u,x\rangle>t\},
\qquad u\in S^{d-1},\ t\in\R,
\]
with a fixed open/closed boundary convention. The problem is to specify an admissible class and halfspace transform for which equality of these data determines the measure.

A natural framework, not supplied as a hypothesis in the record, is signed Radon measures on \(\R^d\setminus\{0\}\), with finite total variation on sets bounded away from zero and suitable control at infinity. For example, the stronger tail condition
\[
|\mu|(\{x:\|x\|\geq\varepsilon\})<\infty\qquad(\varepsilon>0)
\]
would make \(\mu(H(u,t))\) finite for \(t>0\). With a chosen class \(\mathcal C\) and these admissible halfspaces, the desired injectivity statement would be
\[
\mu,\nu\in\mathcal C,\quad
\mu(H(u,t))=\nu(H(u,t))\ \forall u\in S^{d-1},\ t>0
\quad\Longrightarrow\quad\mu=\nu.
\]
The source does not assert that the example tail condition alone suffices; further growth, integrability, or structural hypotheses are what the question asks to determine.

If measures are instead defined on all of \(\R^d\), halfspaces staying away from the origin cannot detect an atom at zero, so that atom must be excluded, specified, or measured by additional data. Halfspaces meeting an infinite-mass singularity also may not have finite signed values: unrestricted subtraction of infinite positive and negative masses is undefined. A principal-value or distributional transform would require its own definition and cannot be inserted silently. Local finiteness on the punctured space alone also does not control unbounded-halfspace tails.

The exact measure domain, sign and variation conventions, admissible halfspaces, near-zero and infinity bounds, and any regularization are absent. The supplied assessment checked 2026-08-21 calls the broad infinite-signed-measure/Radon-transform extension unresolved in the extracted generality. This is a request for a well-defined injectivity class, not an unqualified assertion that all infinite signed measures are determined by possibly undefined halfspace masses.
\subsection{Short English statement}
Find hypotheses making halfspace data determine signed measures that may diverge near zero. The domain, finite-variation tails, admissible halfspaces, and treatment of an origin atom or undefined infinite signed values must be fixed.
\subsection{Sources}
\begin{itemize}
\item Local UnsolvedMath registry, permanent ID 30000708, catalogue entry, \emph{Cram\'er--Wold Uniqueness for Infinite Signed Measures}.
\item \emph{Recent Developments in Financial and Insurance Mathematics and the Interplay with the Industry} (2007), Oberwolfach Reports, pp. 549--590, DOI \url{https://doi.org/10.4171/owr/2007/10}.
\item The supplied statement names the Cram\'er--Wold theorem as background, without a separate bibliographic entry. Supplied literature triage checked 2026-08-21 cites only the report.
\end{itemize}
\end{document}
